% !TeX root = ../thesis.tex
% ******************************* Appendix A **********************************

\chapter{Reproducibility package}
\label{appendix:repro}
\sloppy

This appendix lets someone outside the team rebuild every official number of the thesis and
run the prototype, using only the code repository and the data folder on Google Drive. The
full instructions are in the file \path{docs/REPRODUCE.md} of the repository
\path{github.com/PLHGNAOH/ViVQA-VLM}.

\section{Components}

\begin{table}[htbp]
\caption{Components of the reproducibility package}
\label{tab:repro-parts}
\centering
\small
\begin{tabular}{p{0.22\textwidth}p{0.68\textwidth}}
\toprule
\textbf{Component} & \textbf{Location} \\
\midrule
Source code & \path{src/} (data, adaptation, eval, prompting, prototype) \\
LoRA adapters & Drive \path{ViVQA-VLM/adapters/<run>/epoch_N/}, about 30~MB each, not in git; SHA-256 hashes recorded in \path{configs/official_models.json} \\
Configuration & \path{configs/qwen_lora.yaml} (model, quantisation, LoRA, training, evaluation); \path{configs/pope_vi_objects.json}; the OCR configuration frozen in \path{experiments/W08_ocr_select_dev300_L4_seed42/} \\
Inference scripts & \texttt{python -m src.prototype.infer} (command line) and \texttt{python -m src.prototype.app} (Gradio interface) \\
Experiment logs & \path{experiments/<run>/}: \path{metrics.json}, \path{predictions.json}, \path{run_signature.json}, \path{train_log.json} \\
Environment & \path{requirements-colab.txt} \\
\bottomrule
\end{tabular}
\end{table}

\section{Environment}

Every run uses Google Colab Pro: an L4 GPU with 24~GB for training, evaluation and the
prototype; a T4 GPU for building the OCR cache. Main libraries: PyTorch 2.11.0 (CUDA 13.0, the
version preinstalled on Colab), Transformers 5.17.0, PEFT 0.21.2, bitsandbytes 0.50.2,
Accelerate 1.15.0, qwen-vl-utils 0.0.14 and Gradio 6.31.0. Two notes on installation:

\begin{itemize}
  \item \textbf{Installing GPU PaddleOCR alongside PyTorch} can overwrite PyTorch's NCCL library
  and break \texttt{import torch}. The notebooks include a four-step procedure: record PyTorch's
  \texttt{nvidia-*} packages, install Paddle, reinstall PyTorch's NCCL, and then check both in
  separate processes.
  \item \textbf{VietOCR} is installed with \texttt{pip install --no-deps vietocr==0.3.13},
  because the PyPI release pins an old Pillow version.
\end{itemize}

Each notebook is opened from GitHub in Colab, set to the GPU type stated at the top of the
notebook, and run in full. The notebook checks the hashes of the data, the OCR cache and the
adapter before using the GPU.

\section{Data}

\begin{itemize}
  \item \textbf{ViVQA}: \texttt{python -m src.data.prepare\_vivqa}, producing 10,799 training,
  1,200 validation and 3,001 test questions with the COCO images.
  \item \textbf{ViTextVQA}: the access-controlled Hugging Face repository
  \path{minhquan6203/ViTextVQA} (whoever reproduces the work must request access), built with
  the notebook \path{00_data_vitextvqa_official.ipynb}; the result has 35,159 training, 5,155
  dev and 10,028 test questions, with the test answers taken from
  \path{ViTextVQA_test_gt.json}.
  \item \textbf{OCR cache}: notebook \path{10_ocr_cache_vitextvqa.ipynb} on a T4; the manifest
  with hashes is in \path{experiments/W10_ocr_cache_fulltrain/ocr_cache_manifest.json}.
\end{itemize}

The rebuilt data files must match the \texttt{data\_file\_sha256} hash recorded for each run in
\path{configs/official_models.json}.

\section{The official models}

Table~\ref{tab:repro-models} lists the official models, the notebooks that reproduce them and
the adapter hashes (first 12 characters of the SHA-256 of the weight file). The original table
is generated automatically from the \path{experiments/} directory with
\texttt{python -m src.eval.repro\_manifest --build}, not typed by hand.

\begin{table}[htbp]
\caption{Official models and reproduction notebooks}
\label{tab:repro-models}
\centering
\footnotesize
\begin{tabular}{p{0.19\textwidth}p{0.24\textwidth}p{0.28\textwidth}p{0.16\textwidth}}
\toprule
\textbf{Model} & \textbf{Official result} & \textbf{Notebook (L4)} & \textbf{Adapter SHA-256} \\
\midrule
R0 · reproduction & TextVQA val: VQA-Acc 79.60 & \path{06_r0_textvqa_repro_qwen25vl} & --- \\
B1 · zero-shot & ViVQA EM 22.09; ViTextVQA EM 29.12 & \path{02_run_baseline_qwen25vl}, \path{08_fulltest_vitextvqa_b1_b2} & --- \\
B2 · QLoRA & ViTextVQA EM 50.26 & \path{04_qlora_qwen25vl}, \path{08_fulltest_vitextvqa_b1_b2} & \texttt{2e7b17d18b32} \\
C3 · QLoRA ViVQA & ViVQA EM 76.64 & \path{16_qlora_vivqa_2epoch_select_qwen25vl} & \texttt{3ceff2f7f961} \\
B3-full · OCR-aware QLoRA & ViTextVQA EM 53.35 & \path{18_qlora_ocr_b3full_vs_b2_qwen25vl} & \texttt{c1eb9c2b4baf} \\
SH · shared (prototype) & ViVQA EM 73.51; ViTextVQA test2000 EM 49.50 & \path{20_qlora_shared_adapter_qwen25vl} & \texttt{c9791e961aaa} \\
\bottomrule
\end{tabular}
\end{table}

\section{Commands}

Every notebook calls the two basic commands below with preset arguments. Example for B3-full:

\begin{small}
\begin{verbatim}
python -m src.adaptation.train_qlora --dataset vitextvqa_official \
    --run_name <run> --epochs 1 --save_steps 100 \
    --mode ocr --ocr_backend paddledet_vietocr --ocr_template v1 \
    --ocr_max_chars 600 --ocr_min_conf 0.0

python -m src.eval.run_baseline --dataset vitextvqa_official \
    --data_file <.../test.json> --run_name <run_eval> --chunk 500 \
    --adapter <Drive>/adapters/<run>/epoch_1 \
    --mode ocr --ocr_backend paddledet_vietocr --ocr_template v1 \
    --ocr_max_chars 600 --ocr_min_conf 0.0
\end{verbatim}
\end{small}

Removing the OCR arguments (\texttt{--mode ocr} and \texttt{--ocr\_*}) gives B2 or C3; removing
\texttt{--adapter} gives B1. The shared configuration (bf16, 4-bit NF4, at most 401,408 pixels,
greedy decoding, at most 32 generated tokens, seed 42, LoRA $r = 16$ on q, k, v, o) is in
\path{configs/qwen_lora.yaml}. The file \path{run_signature.json} prevents resuming an
unfinished run with a different configuration.

The analyses (full evaluation table, paired tests, three seeds, the $2\times2$ experiment, LoRA
ablations, POPE-vi, error analysis, the EXIF check) need only a CPU and re-read
\path{predictions.json}; the corresponding commands and notebooks are listed in
\path{docs/REPRODUCE.md}, section 5.

\section{Prototype}

The interface runs with the notebook \path{28_prototype_v2_live_demo_sh_qwen25vl.ipynb} on an
L4; the notebook re-checks rules R-OCR and R-LIVE and then opens the Gradio link. It can also
be run from the command line:

\begin{small}
\begin{verbatim}
python -m src.prototype.infer --official sh \
    --drive_root /content/drive/MyDrive/ViVQA-VLM \
    --image image.jpg --question "<question>" --out result.jsonl
\end{verbatim}
\end{small}

The argument \texttt{--official sh} picks the SH adapter and checks its hash automatically;
adding \texttt{--mock} allows a test without a GPU.

\section{Determinism and known limitations}

Evaluation uses greedy decoding and runs question by question. Re-running on the same GPU type
(L4) gives 100\% identical predictions (e.g.\ 9,902/9,902 questions in the week-14 check).
Changing the GPU type (T4 versus L4) can shift the OCR slightly on a few images, so the GPU type
stated in the notebook should be used. Known limitations:

\begin{enumerate}
  \item Some libraries (bitsandbytes, Accelerate, qwen-vl-utils) are not recorded in each run's
  log but only in the environment snapshot; the two records agree.
  \item PyTorch is the version preinstalled on Colab; if Colab changes the default, the version
  recorded in the logs must be installed.
  \item The official OCR cache has the EXIF bug for 1.256\% of the ViTextVQA test questions; the
  official numbers are kept, and the post-fix results are reported as a sensitivity analysis
  (Section~\ref{sec:res-exif}). To reproduce the official cache exactly, use the source code
  before the fix (commit \texttt{b9f7fd4}).
  \item ViTextVQA is access-controlled: whoever reproduces the work must request access.
\end{enumerate}
